
%% ============================================================
%%  Akronyme
%% ============================================================

\newacronym{api}{API}{Application Programming Interface}
\newacronym{rest}{REST}{Representational State Transfer}
\newacronym{http}{HTTP}{Hypertext Transfer Protocol}
\newacronym{ocr}{OCR}{Optical Character Recognition}
\newacronym{nlp}{NLP}{Natural Language Processing}
\newacronym{cnn}{CNN}{Convolutional Neural Network}
\newacronym{tfidf}{TF-IDF}{Term Frequency -- Inverse Document Frequency}
\newacronym{saas}{SaaS}{Software as a Service}
\newacronym{rbac}{RBAC}{Role-Based Access Control}
\newacronym{abac}{ABAC}{Attribute-Based Access Control}
\newacronym{nist}{NIST}{National Institute of Standards and Technology}
\newacronym{dto}{DTO}{Data Transfer Object}
\newacronym{acl}{ACL}{Access Control List}
\newacronym{jwt}{JWT}{JSON Web Token}
\newacronym{pdf}{PDF}{Portable Document Format}
\newacronym{sql}{SQL}{Structured Query Language}
\newacronym{sha}{SHA-256}{Secure Hash Algorithm 256}
\newacronym{css}{CSS}{Cascading Style Sheets}
\newacronym{vit}{ViT}{Vision Transformer}
\newacronym{vlm}{VLM}{Vision-Language Model}
\newacronym{llm}{LLM}{Large Language Model}
\newacronym{crm}{CRM}{Customer Relationship Management}
\newacronym{odata}{OData}{Open Data Protocol}
\newacronym{mime}{MIME}{Multipurpose Internet Mail Extensions}

%% ============================================================
%%  Glossar
%% ============================================================

\newglossaryentry{ki}
{
    name=Künstliche Intelligenz,
    description={Teilgebiet der Informatik, das sich mit der Entwicklung von Systemen befasst, die Aufgaben ausführen können, die üblicherweise menschliche Intelligenz erfordern, etwa Klassifikation, Sprachverarbeitung oder Mustererkennung}
}

\newglossaryentry{dokumentenklassifikation}
{
    name=Dokumentenklassifikation,
    description={Automatisierte Zuordnung eingehender Dokumente zu vordefinierten Klassen anhand von Text-, Layout- und visuellen Merkmalen}
}

\newglossaryentry{embedding}
{
    name=Embedding,
    description={Dichte Vektorrepräsentation eines Objekts in einem kontinuierlichen Raum, in dem semantisch ähnliche Objekte nahe beieinanderliegen}
}

\newglossaryentry{merkmalsextraktion}
{
    name=Merkmalsextraktion,
    description={Verfahren zur Umwandlung von Rohdaten (Text, Bilder, Layout) in numerische oder strukturierte Merkmale, die von Lernverfahren verarbeitet werden können}
}

\newglossaryentry{microservices}
{
    name=Microservices,
    description={Architekturansatz, bei dem eine Anwendung in kleine, unabhängig deploybare Dienste aufgeteilt wird, die über definierte Schnittstellen kommunizieren}
}

\newglossaryentry{monolith}
{
    name=Monolithische Architektur,
    description={Klassischer Architekturansatz, bei dem alle Funktionsbereiche einer Anwendung in einer einzigen, eng gekoppelten Einheit zusammengefasst und gemeinsam deployt werden}
}

\newglossaryentry{blazor}
{
    name=Blazor WebAssembly,
    description={Microsoft-Framework zur Entwicklung interaktiver Webanwendungen mit C\#, wobei der Anwendungscode clientseitig im Browser über WebAssembly ausgeführt wird}
}

\newglossaryentry{postgresql}
{
    name=PostgreSQL,
    description={Quelloffenes, relationales Datenbankmanagementsystem mit umfangreichen SQL-Funktionalitäten und hoher Skalierbarkeit}
}

\newglossaryentry{azureaisearch}
{
    name=Azure AI Search,
    description={Cloud-basierter Suchdienst von Microsoft zur Indizierung großer Dokumentbestände mit Volltext-, Filter- und Facettensuche}
}

\newglossaryentry{entraid}
{
    name=Microsoft Entra ID,
    description={Cloudbasierter Identitäts- und Zugriffsverwaltungsdienst von Microsoft für Authentifizierung, Autorisierung und Schutz von Benutzern, Geräten und Anwendungen}
}

\newglossaryentry{sharepoint}
{
    name=SharePoint Online,
    description={Cloudbasierte Kollaborations- und Dokumentenmanagementplattform von Microsoft, bereitgestellt als SaaS im Rahmen von Microsoft~365}
}

\newglossaryentry{circuitbreaker}
{
    name=Circuit Breaker,
    description={Entwurfsmuster, das bei anhaltenden Fehlern einer Abhängigkeit weitere Aufrufe unterbindet und so eine sanfte Degradation des Gesamtsystems ermöglicht}
}

\newglossaryentry{retrypattern}
{
    name=Retry Pattern,
    description={Entwurfsmuster, bei dem fehlgeschlagene Aufrufe kontrolliert und mit steigender Verzögerung wiederholt werden, um vorübergehende Störungen zu überbrücken}
}

\newglossaryentry{gateway}
{
    name=API-Gateway,
    description={Vorgeschaltete Vermittlungsschicht, die Anfragen an Backend-Dienste weiterleitet und zentrale Aufgaben wie Authentifizierung, Routing, Quotensteuerung und Monitoring übernimmt}
}

\newglossaryentry{zerotrust}
{
    name=Zero Trust,
    description={Sicherheitsstrategie, bei der Identitäten, Zugriffsbedingungen und Berechtigungen konsequent überprüft werden, unabhängig davon, ob sich der Zugriff innerhalb oder außerhalb des Netzwerks befindet}
}

\newglossaryentry{tokenization}
{
    name=Tokenisierung,
    description={Zerlegung von Text in einzelne Einheiten (Tokens), die als Eingabe für Sprachmodelle oder Textverarbeitungsverfahren dienen}
}

\newglossaryentry{transformer}
{
    name=Transformer,
    description={Neuronale Netzwerkarchitektur, die auf Selbstaufmerksamkeit (Self-Attention) basiert und Grundlage moderner Sprachmodelle wie BERT oder GPT bildet}
}

\newglossaryentry{cosinesimilarity}
{
    name=Kosinus-Ähnlichkeit,
    description={Maß für die Ähnlichkeit zweier Vektoren, berechnet als Kosinus des Winkels zwischen ihnen; wird häufig zum Vergleich von Embeddings eingesetzt}
}

\newglossaryentry{preflight}
{
    name=Preflight,
    description={Vorverarbeitungsphase im Upload-Prozess, in der ein Dokument vor der finalen Persistierung durch die KI-Klassifikation analysiert und die Ergebnisse dem Benutzer zur Bestätigung vorgelegt werden}
}

\newglossaryentry{facet}
{
    name=Facette,
    description={Dynamisch berechnete Filteroptionen in Suchergebnissen, die auf den tatsächlich vorhandenen Werten eines Attributs basieren und die Eingrenzung der Treffermenge ermöglichen}
}
